% TOP_PROBLEM: {"id":30007648,"problem_number":"LOCAL-30007648","title":"Social connections as a causal mobility tool","category_id":21,"category":{"id":21,"name":"economics","display_name":"Economics","description":"Open economic theory statements, published research agendas, and proposed scoped specifications; question provenance and historical priority are qualified per record.","slug":"economics","order_index":21,"created_at":"2026-10-08T00:00:00Z"},"status":"research_question","external_url":null,"legacy_ids":[],"published":false,"statement":"In the explicitly specified setting: Does deliberately increasing cross-class friendships raise disadvantaged children's adult incomes, and through which information, aspiration, referral, or support mechanisms? The mathematical statement above fixes the target, assumptions and scope of this catalogue formulation.","economics":{"original_id":137,"field":"Inequality and economic mobility","kind":"identification","formalization_basis":"adapted_research_question","source_status":"explicit_open","priority_status":"earliest_found","priority_note":"Earliest directly inspected qualifying passage in this audit; earlier origins were not established. A review or research-agenda statement does not imply original priority.","earliest_verified_reference":{"authors":["Raj Chetty","Matthew O. Jackson","Theresa Kuchler","Johannes Stroebel","Nathaniel Hendren","Robert B. Fluegge","Sara Gong","Federico Gonzalez","Armelle Grondin","Matthew Jacob","Drew Johnston","Martin Koenen","Eduardo Laguna-Muggenburg","Florian Mudekereza","Tom Rutter","Nicolaj Thor","Wilbur Townsend","Ruby Zhang","Mike Bailey","Pablo Barberá","Monica Bhole","Nils Wernerfelt"],"title":"Social capital I: measurement and associations with economic mobility","year":2022,"url":"https://tomrutter42.github.io/assets/pdf/sc1.pdf","doi":"10.1038/s41586-022-04996-4","locator":"Discussion, Nature printed p. 120 (PDF p. 13), final research-agenda paragraph","evidence_summary":"The authors explicitly request direct study of whether increasing economic connectedness raises intergenerational income mobility."},"current_reference":{"authors":["Raj Chetty","Matthew O. Jackson","Theresa Kuchler","Johannes Stroebel","Nathaniel Hendren","Robert B. Fluegge","Sara Gong","Federico Gonzalez","Armelle Grondin","Matthew Jacob","Drew Johnston","Martin Koenen","Eduardo Laguna-Muggenburg","Florian Mudekereza","Tom Rutter","Nicolaj Thor","Wilbur Townsend","Ruby Zhang","Mike Bailey","Pablo Barberá","Monica Bhole","Nils Wernerfelt"],"title":"Social capital I: measurement and associations with economic mobility","year":2022,"url":"https://tomrutter42.github.io/assets/pdf/sc1.pdf","doi":"10.1038/s41586-022-04996-4","locator":"Discussion, Nature printed p. 120 (PDF p. 13), final research-agenda paragraph","evidence_summary":"The authors explicitly request direct study of whether increasing economic connectedness raises intergenerational income mobility."},"formalization_note":"The cited passage establishes a research direction, not this exact mathematical statement. The notation, population, policy menu, assumptions, and estimands are an original adaptation for this catalogue; no universal unsolved theorem or source priority is claimed."},"statusQualification":"Published question or research agenda verified at the cited locator. The mathematical specification is adapted for this catalogue and includes additional assumptions; source publication does not establish first-ever priority or certify this exact model remains unsolved. The cited passage establishes a research direction, not this exact mathematical statement. The notation, population, policy menu, assumptions, and estimands are an original adaptation for this catalogue; no universal unsolved theorem or source priority is claimed.","statusReviewedAt":"2026-10-08"}
% FIRST_POSTED: 2026-10-08
% ENTRY_KIND: statement_only
% !TeX program = lualatex
\documentclass[11pt]{article}
\usepackage{fontspec}
\usepackage[a4paper,margin=25mm]{geometry}
\usepackage{amsmath,amssymb,mathtools}
\usepackage{xurl}
\usepackage[hidelinks]{hyperref}
\newcommand{\sref}[2]{\hyperlink{source:#1:#2}{[S#2]}}
\setlength{\emergencystretch}{3em}
\begin{document}
% problemId:30007648
\section*{Social connections as a causal mobility tool}\label{30007648}
\subsection{Definitions and mathematical statement}
Fix a cohort of children from low-income families and a network-forming intervention \(z\in\{0,1\}\). The intervention is precisely specified, such as a common activity with randomized cross-income participation; it does not directly give children higher income. Let \(A_i(z)\) be the share of child \(i\)'s friends from a predeclared higher socioeconomic group, and let \(Y_i(z)\) be adult income rank relative to a fixed reference distribution. Define the total policy effect \(\tau=E[Y_i(1)-Y_i(0)]\) and the network first stage \(\phi=E[A_i(1)-A_i(0)]\). Estimate both and determine whether increased connectedness improves mobility. When \(\phi\ne0\), interpreting \(\tau/\phi\) as an effect of the continuous friendship share requires explicit exclusion, interference restrictions, and either a justified local-average-response model or a linear response restriction; exposure or program participation may affect outcomes through other routes. Measure candidate mechanisms including information, aspirations, mentorship, and referrals with separately justified designs. Account for friendship selection, persistence of connections, attrition, and displacement of opportunities for untreated children. The task adapts the authors' explicit call for direct tests of connectedness interventions. A strong area-level correlation, including correlation with causal place effects, does not by itself identify the causal effect of creating these friendships.

\par\textbf{Formulation and scope.} The cited passage establishes a research direction, not this exact mathematical statement. The notation, population, policy menu, assumptions, and estimands are an original adaptation for this catalogue; no universal unsolved theorem or source priority is claimed.

\par\textbf{Open-question provenance.} An explicit question or research agenda was verified at the source locator. This does not by itself certify that every later version remains unresolved \sref{30007648}{1}.

\par\textbf{Historical priority.} Earliest directly inspected qualifying passage in this audit; earlier origins were not established. A review or research-agenda statement does not imply original priority.

\subsection{Short English statement}
In the explicitly specified setting: Does deliberately increasing cross-class friendships raise disadvantaged children's adult incomes, and through which information, aspiration, referral, or support mechanisms? The mathematical statement above fixes the target, assumptions and scope of this catalogue formulation.

\subsection{Sources}
\begin{itemize}\raggedright
\item[\textnormal{[S1]}]\hypertarget{source:30007648:1}{} Raj Chetty, Matthew O. Jackson, Theresa Kuchler, Johannes Stroebel, Nathaniel Hendren, Robert B. Fluegge, Sara Gong, Federico Gonzalez, Armelle Grondin, Matthew Jacob, Drew Johnston, Martin Koenen, Eduardo Laguna-Muggenburg, Florian Mudekereza, Tom Rutter, Nicolaj Thor, Wilbur Townsend, Ruby Zhang, Mike Bailey, Pablo Barberá, Monica Bhole, Nils Wernerfelt. 2022. Social capital I: measurement and associations with economic mobility. Discussion, Nature printed p. 120 (PDF p. 13), final research-agenda paragraph.
\url{https://tomrutter42.github.io/assets/pdf/sc1.pdf}.
DOI: 10.1038/s41586-022-04996-4.
{\small\textit{Earliest verified question/agenda reference; historical priority qualified below. The authors explicitly request direct study of whether increasing economic connectedness raises intergenerational income mobility.}}
\end{itemize}
\end{document}
